% =============================================================
%  CONCLUSIONES
% =============================================================
\renewcommand{\seccorta}{Conclusiones}
\section{Conclusiones}

\begin{frame}{Conclusiones}
  \begin{columns}[T]
    \column{0.48\textwidth}
    \begin{block}{Enlace A}
      \small
      \begin{itemize}
        \item \EAunoAlt/\EAdosAlt\ son las alturas mínimas para salvar la edificación
              con $\geq$60\,\% de Fresnel.
        \item $P_{RX}$ teórica \EAprxTeo\ coincide con la simulada \EAprxSim.
        \item Canal de 20 MHz: SNR \EAsnr\ y $\sim$63 Mbps útiles $>$ 40 Mbps requeridos.
        \item Se equilibran SNR, capacidad y normativa sin máxima potencia ni altura.
      \end{itemize}
    \end{block}
    \column{0.48\textwidth}
    \begin{block}{Enlace B}
      \small
      \begin{itemize}
        \item La validación 3D en Google Earth fue indispensable porque UISP subestima los edificios.
        \item Con \EBunoAlt/\EBdosAlt\ se despeja el \EBdespejePct\ de Fresnel ($A_{obs}=0$).
        \item $P_{RX}$ teórica \EBprxTeo\ vs.\ simulada \EBprxSim: diferencia $\approx$1.2 dB.
        \item Con ruido urbano: SNR de \EBsnr\ $\Rightarrow$ 16QAM; $\approx$64 Mbps útiles $>$ 50 Mbps requeridos.
      \end{itemize}
    \end{block}
  \end{columns}
  \vspace{4pt}
  \begin{alertblock}{Generales}
    \small En entornos urbanos sin LIDAR, el obstáculo real son los edificios y debe validarse en 3D.
    El piso de ruido supuesto define la modulación y la capacidad: el diseño es un compromiso entre capacidad, robustez frente al ruido y costo; la modulación
    adaptativa mantiene el enlace ante degradaciones, y en banda compartida la potencia debe
    ser la mínima necesaria dentro del límite de PIRE.
  \end{alertblock}
\end{frame}

\begin{frame}[plain]
  \centering\vfill
  {\usebeamercolor[fg]{title}\Huge\bfseries ¡Gracias!}\\[8pt]
  {\color{naranja}\large ¿Preguntas?}
  \vfill
\end{frame}
